\documentclass{article}
\usepackage{amsmath, amssymb, amsthm}

\newtheorem{theorem}{Theorem}

\begin{document}

\section*{Multiple Testing and FCR Control}

\subsection*{Setup}

For each hypothesis $i = 1,\dots,m$, suppose we observe data
\[
Z_i \sim \mathcal{N}(\mu_i, \sigma_i^2),
\]
where $\sigma_i^2$ is known (or consistently estimated).

We test
\[
H_{0,i}: \mu_i = 0.
\]

The usual two-sided $p$-value is
\[
P_i = 2\left(1 - \Phi\left(\frac{|Z_i|}{\sigma_i}\right)\right),
\]
where $\Phi$ is the standard normal CDF.

A $(1-\alpha)$ confidence interval for $\mu_i$ is
\[
C_i(1-\alpha)
=
\left[
Z_i - z_{1-\alpha/2}\sigma_i,
\;
Z_i + z_{1-\alpha/2}\sigma_i
\right].
\]

\subsection*{Naive Procedure}

\begin{enumerate}
\item Apply the Benjamini–Hochberg (BH) procedure at level $q$ to $P_1,\dots,P_m$.
\item Let $I \subset \{1,\dots,m\}$ be the set of discoveries.
\item Report $C_i(1-\alpha)$ for each $i \in I$.
\end{enumerate}

\textbf{Problem:} This does not control the \emph{False Coverage Rate} (FCR), because intervals are constructed after selection.

\subsection*{False Coverage Rate (FCR)}

Let
\[
V = \#\{ i \in I : \mu_i \notin C_i \},
\qquad
R = |I|.
\]

Define the False Coverage Proportion (FCP) as
\[
\mathrm{FCP} =
\begin{cases}
\frac{V}{R}, & R > 0, \\
0, & R = 0.
\end{cases}
\]

The False Coverage Rate is
\[
\mathrm{FCR} = \mathbb{E}[\mathrm{FCP}]
=
\mathbb{E}\left[
\frac{\#\{ i \in I : \mu_i \notin C_i \}}{\max(R,1)}
\right].
\]

Goal: construct a procedure such that $\mathrm{FCR} \le \alpha$.

\subsection*{FCR-Controlling Procedure}

\begin{enumerate}
\item Apply BH at level $q$ to obtain selected set $I$.
\item Let $R = |I|$.
\item For each $i \in I$, construct a confidence interval at level
\[
1 - \frac{qR}{m}.
\]
\end{enumerate}

Thus the reported interval is
\[
C_i\left(1 - \frac{qR}{m}\right).
\]

\subsection*{Theorem (FCR Control)}

\begin{theorem}
Assume:
\begin{enumerate}
\item The data across $i=1,\dots,m$ are independent.
\item For each $i$, $C_i(1-\alpha)$ is a valid $(1-\alpha)$ confidence interval for $\mu_i$.
\item The intervals are constructed from data used in test $i$.
\end{enumerate}
Then the above procedure satisfies
\[
\mathrm{FCR} \le q.
\]
\end{theorem}

\subsection*{Key Idea of Proof}

Write
\[
\mathrm{FCR}
=
\mathbb{E}
\left[
\sum_{i=1}^m
\frac{\mathbf{1}\{ i \in I, \mu_i \notin C_i(1 - qR/m)\}}
{\max(R,1)}
\right].
\]

Condition on all $p$-values except $P_i$.
Under independence, the expected contribution from each $i$ is bounded by
\[
\frac{q}{m}.
\]

Summing over $i=1,\dots,m$ yields
\[
\mathrm{FCR} \le q.
\]

\subsection*{More General Framework}

Suppose for each $i$ we have a nested family of valid confidence intervals
\[
C_i(1-\alpha_1) \subseteq C_i(1-\alpha_2)
\quad \text{for } \alpha_1 < \alpha_2.
\]

Define
\[
P_i = \inf\{\alpha : 0 \notin C_i(1-\alpha)\}.
\]

Then applying BH to $P_1,\dots,P_m$ and reporting
\[
C_i\left(1 - \frac{qR}{m}\right)
\]
for selected $i$ controls FCR at level $q$.

\subsection*{Connection to BH}

The BH procedure selects the largest $k$ such that
\[
P_{(k)} \le \frac{qk}{m}.
\]

Equivalently, we want as many $p$-values as possible below their BH thresholds.

After selection, confidence intervals are adjusted to match the BH threshold $\frac{qR}{m}$.

\subsection*{Summary}

Two main perspectives:

\begin{enumerate}
\item Given hypotheses $H_{0,i}$ and valid $p$-values $P_i$, control FDR using BH.
\item Given parameters $\mu_i$ and valid confidence intervals $C_i(1-\alpha)$, adjust coverage after selection to control FCR.
\end{enumerate}

\subsection*{Additional Questions}

\begin{itemize}
\item How to incorporate prior information?
\item How to handle streaming or online data?
\item How to adjust for dependence across tests?
\end{itemize}

\end{document}
